% Reconstructed text from V5 PDF page 3; non-native source
\par\smallskip\noindent\rule{\textwidth}{0.4pt}\par{\sffamily\footnotesize Original V5 page 3 -- recovered text}\par\begin{flushleft}\footnotesize\ttfamily
\par\smallskip
\par\smallskip
\noindent Diagnostics and Serving Infrastructure for MAS 3\par
\par\smallskip
\par\smallskip
\noindent Table 1: Conceptual differences from adjacent reviews. The entries describe emphases rather than binary\par
\noindent claims of exclusive coverage. Our contribution is a comparison of evidence requirements and state-\par
\noindent lifetime dependencies across diagnosis and serving, not a claim to invent the component techniques.\par
\par\smallskip
\noindent  Review area Unit of analysis Principal concern Dimension added here\par
\noindent  MAS orchestration Agent roles, workflow Coordination, framework Serving-state lifetimes and diagnostic\par
\noindent  graph design dependencies\par
\noindent  AgentOps Telemetry, operational Monitoring, trace collec- Evidence preservation vs. reclamation\par
\noindent  practice tion cost\par
\noindent  Agent evaluation Benchmark tasks, scores Offline task quality Replayability, granularity, online timing\par
\noindent  LLM-serving Requests, KV blocks, Latency, throughput Agent-level future calls and restoration\par
\noindent  workers\par
\noindent  MAS collaboration Inter-agent interactions Failure, collaboration pat- Shared evidence, explicit retention con-\par
\noindent  terns ditions\par
\noindent  This review Diagnostic evidence, Conditional cross-layer Evidence inputs, training, replay, tim-\par
\noindent  serving state dependencies ing, cost, recoverability\par
\par\smallskip
\par\smallskip
\par\smallskip
\noindent 1.2 The Diagnostic-Infrastructure Gap\par
\noindent The gap between diagnostic and infrastructure research is both conceptual and practical. Con-\par
\noindent ceptually, diagnostic methods are designed under an implicit assumption that execution traces\par
\noindent are always available — stored in a log database, replay buffer, or trace store that is independent\par
\noindent of the serving engine’s memory management. Infrastructure systems, conversely, optimize for\par
\noindent throughput and latency without considering whether their eviction decisions affect downstream\par
\noindent diagnostic capability. Neither community has systematically characterized the conditions under\par
\noindent which these assumptions interact.\par
\noindent  Practically, this gap manifests in three ways. First, among the 12 evaluation benchmarks\par
\noindent in our curated corpus, we did not identify one that evaluates how diagnostic overhead affects\par
\noindent serving throughput or how eviction policies affect attribution accuracy. Second, within the\par
\noindent serving-infrastructure studies examined at full-text depth, we did not identify a reported sys-\par
\noindent tem that jointly optimizes diagnostic retention and serving throughput under a single measured\par
\noindent resource budget — the closest system (TOPAS) optimizes serving throughput alone. Third,\par
\noindent the observability standards that could bridge the gap (OpenTelemetry GenAI conventions) lack\par
\noindent serving-side attributes: there are no standardized attributes for KV cache residency, eviction\par
\noindent events, or memory pressure.\par
\noindent  This review addresses the gap by mapping both layers along shared dimensions (evidence\par
\noindent inputs, training, replay, granularity, timing, cost), formalizing the conditions for cross-layer\par
\noindent conflict (Section 6), and proposing a structured evaluation framework (DRICE, Section 8) and\par
\noindent validation protocol (Section 7.3) that make the shared-budget hypothesis testable.\par
\noindent  This is a curated critical review rather than a systematic review. The inclusion boundary,\par
\noindent evidence-selection rules, and verification tiers are described in Section 3. Quantitative claims\par
\noindent used for comparison come from full-text sources; abstract-level results are marked [a] and are\par
\noindent not used to compute cross-system deltas or rankings.\par
\par\smallskip
\noindent Why gap identification is a contribution. Identifying that a widely assumed conflict is con-\par
\noindent ditional (not universal) changes the design space for joint diagnostic-infrastructure systems. If\par
\noindent the conflict were universal, practitioners would need to build shared-budget managers for every\par
\noindent deployment. Because the conflict requires three specific conditions (Eq. 1), practitioners should\par
\noindent instead instrument for the specific conditions under which interference occurs. The simulated\par
\noindent KV-eviction experiment (Section 6.6) illustrates that, in a simplified CPU-based cache model,\par
\noindent naive policy augmentation produces a negative trade-off; live GPU-serving validation remains\par
\noindent future work. The retention ablation (Section 7.3) reveals a further non-obvious interaction on\par
\end{flushleft}
